\chapter{Axiome der IWT}

\section{Einleitung}

Die Informations-Weber-Theorie (IWT) ist eine fundamentale, diskrete Ur-Theorie. Sie postuliert keine Raumzeit, keine Materie, keine Felder – nur Information. Raum, Zeit, Materie und Energie emergieren aus der Struktur und Dynamik eines diskreten Informationsnetzwerks.

Die folgenden neun Axiome definieren die IWT vollständig. Sie sind logisch unabhängig und bilden die Grundlage für alle weiteren Ableitungen.

\section{Axiom 1: Diskretes universelles Informationsfeld}

Es existiert ein skalares, diskretes Informationsfeld \( I_k^{(n)} \), das die vollständige physikalische Realität beschreibt.

\[
I_k^{(n)} \in \mathbb{C}
\]

Hierbei ist:
\begin{itemize}
\item \( k = 1, \dots, N \) – Index der diskreten Informationsknoten,
\item \( n \in \mathbb{N}_0 \) – diskreter Zeitindex (Update-Schritt),
\item \( I_k^{(n)} \) – komplexer Informationswert.
\end{itemize}

Die Amplitude \( |I_k^{(n)}| \) kodiert die Stärke der Information am Knoten. Die Phase \( \arg(I_k^{(n)}) \) kodiert die kohärenten Relationen im Netzwerk.

Materie, Energie, Wellen, Geometrie und Dynamik sind Manifestationen der Struktur und Veränderung dieses komplexen Feldes.

\section{Axiom 2: Informationserhaltung}

Die Gesamtinformation – definiert als die Summe der Betragsquadrate – ist erhalten:

\[
\sum_{k=1}^{N} |I_k^{(n)}|^2 = \text{const} \quad \text{für alle } n
\]

Information wird nicht erzeugt oder vernichtet – sie wird nur zwischen den Knoten umverteilt.

\section{Axiom 3: Informationsmetrik als emergente Geometrie}

Die physikalische Raumzeit ist nicht fundamental. Stattdessen entsteht eine effektive diskrete Metrik \( g_{kl}^{(n)} \) aus der lokalen und globalen Struktur des Informationsfeldes.

Die Kopplungsmatrix \( K_{kl}^{(n)} \) beschreibt die Stärke der direkten Wechselwirkung zwischen zwei Knoten. Sie emergiert aus der fraktalen Geometrie des Raumes:

\[
K_{kl}^{(n)} = \frac{1}{d_{kl}^{3-D}}
\]

Dabei ist \( d_{kl} \) die fraktale Distanz zwischen den Knoten \( k \) und \( l \), und \( D \) die fraktale Dimension des Raumes.

Die diskrete Informationsmetrik \( g_{kl}^{(n)} \) ist der normalisierte Zusammenhang dieser Kopplungsmatrix:

\[
g_{kl}^{(n)} = \frac{K_{kl}^{(n)}}{\sqrt{K_{kk}^{(n)} K_{ll}^{(n)}}}
\]

\section{Axiom 4: Variationsprinzip der diskreten Informationsdynamik}

Die Dynamik des Informationsfeldes folgt aus einem universellen diskreten Informations-Lagrange-Funktional:

\[
\mathcal{L}_d[I_k^{(n)}] = \mathcal{L}_{\text{local}} + \mathcal{L}_{\text{global}} + \mathcal{L}_{\text{fraktal}}
\]

Die Variation des Funktionals liefert die Bewegungsgleichungen:

\[
\delta \mathcal{L}_d = 0
\]

Daraus folgt die diskrete Euler-Lagrange-Gleichung:

\[
\frac{\partial \mathcal{L}_d}{\partial I_k^{(n)}} - \Delta_t^+ \left( \frac{\partial \mathcal{L}_d}{\partial (\Delta_t I_k^{(n)})} \right) - \Delta_s \left( \frac{\partial \mathcal{L}_d}{\partial (\Delta_s I_k^{(n)})} \right) = 0
\]

\section{Axiom 5: Dynamische Gleichung der Informationsmetrik}

Die Metrik entsteht aus der Variation des diskreten Funktionals. Die fundamentale Gleichung der Informationsmetrik in diskreter Form lautet:

\[
g_{kl}^{(n+1)} = g_{kl}^{(n)} + T \left[ \Delta_k I^{(n)} \Delta_l I^{(n)} - \lambda \frac{\Delta_{kl}^2 I^{(n)}}{I_{kl}^{(n)}} + \mu g_{kl}^{(n)} \ln(1 + \gamma_{\text{eff}} G_{\text{eff}} L^2) \right]
\]

Sie vereinigt lokale Weber-Dynamik, globale Bohm-Struktur und fraktale kosmische Skalierung.

\section{Axiom 6: Energie als emergenter Informationsfluss}

Energie ist keine fundamentale Größe. Sie entsteht als Erhaltungsgröße des diskreten Informationsflusses.

Die diskrete Energie folgt aus der Zeitinvarianz des Informations-Lagrange-Funktionals:

\[
E^{(n)} = \sum_k \left[ \frac{\partial \mathcal{L}_d}{\partial (\Delta_t I_k^{(n)})} \Delta_t I_k^{(n)} - \mathcal{L}_d \right]
\]

Und es gilt:

\[
E^{(n+1)} = E^{(n)}
\]

\section{Axiom 7: Naturkonstanten als emergente Skalierungsparameter}

Die fundamentalen Naturkonstanten sind keine Eingaben der Theorie, sondern feste Punkte der diskreten Informationsdynamik. Sie emergieren aus den Netzwerkparametern:

\begin{align*}
c &= \frac{\lambda_0}{T} \\
\hbar &= \alpha_{\text{IWT}} \cdot \Delta I_{\text{min}} \cdot \lambda_0^2 \\
G &= \beta_{\text{IWT}} \cdot \frac{\lambda_0^{3-D}}{T^2} \\
\alpha &= \frac{\alpha_{\text{IWT}}}{\beta_{\text{IWT}}}
\end{align*}

\section{Axiom 8: Emergenz von Raum, Zeit und Dynamik}

Raum, Zeit und Dynamik sind emergente Eigenschaften der Informationsmetrik:

\begin{itemize}
\item \textbf{Raum:} Der physikalische Raum entsteht aus der diskreten Metrik \( g_{kl}^{(n)} \), die selbst aus der fraktalen Kopplungsmatrix folgt.

\[
ds_{kl}^{(n)} = \sqrt{g_{kl}^{(n)}} \cdot \lambda_0
\]

\item \textbf{Zeit:} Die physikalische Zeit entsteht als Ordnungsstruktur der Informationsänderung:

\[
t \approx n \cdot T
\]

wobei \( n \) der Update-Index und \( T \) der fundamentale Zeitschrift ist.

\item \textbf{Dynamik:} Klassische Mechanik, Quantenmechanik und Gravitation sind Grenzfälle der diskreten Informationsdynamik.
\end{itemize}

\section{Axiom 9: Universelle Gültigkeit}

Die IWT gilt auf allen Skalen:

\begin{itemize}
\item \textbf{Mikroskopisch:} Quantenstruktur, Elementarteilchen, Wellenphänomene.
\item \textbf{Mesoskopisch:} Klassische Mechanik, Elektrodynamik, Thermodynamik.
\item \textbf{Makroskopisch:} Gravitation, Astrophysik, Planetenbewegung.
\item \textbf{Kosmologisch:} Rotverschiebung, Hintergrundstrahlung, großskalige Struktur.
\end{itemize}

\section{Zusammenfassung der Axiome}

Die neun Axiome definieren die IWT vollständig in ihrer fundamentalen diskreten Formulierung.

\begin{enumerate}
\item Diskretes universelles Informationsfeld \( I_k^{(n)} \)
\item Informationserhaltung \( \sum |I|^2 = \text{const} \)
\item Informationsmetrik \( g_{kl}^{(n)} \) aus Kopplungsmatrix
\item Variationsprinzip \( \delta \mathcal{L}_d = 0 \)
\item Dynamische Gleichung der Informationsmetrik
\item Energie als emergenter Informationsfluss
\item Naturkonstanten als emergente Skalierungsparameter
\item Emergenz von Raum, Zeit und Dynamik
\item Universelle Gültigkeit
\end{enumerate}

Alle physikalischen Phänomene – von der Quantenmechanik über die Gravitation bis zur Kosmologie – folgen aus der Struktur und Dynamik des diskreten Informationsfeldes und der daraus emergierenden Metrik.

\section{Q ist kein Axiom}

\( Q \) ist nicht in den Axiomen enthalten.

\( Q \) ist kein Postulat. \( Q \) ist ein Resultat.

\( Q \) emergiert aus der Selbstorganisation der Information, die durch die Axiome definiert ist.

Die Herleitung von \( Q \) aus den Axiomen wird in Kapitel 3 gezeigt.